% manuscript.tex — Adhesion-controlled rigidity transition in a multi-cell phase field model
% Target: Physical Review E  (revtex4-2, two-column)
%
% Build:  pdflatex manuscript && bibtex manuscript && pdflatex manuscript && pdflatex manuscript
%
\documentclass[aps,pre,twocolumn,superscriptaddress,longbibliography]{revtex4-2}

% ── packages ──────────────────────────────────────────────────────────────────
\usepackage{amsmath,amssymb}
\usepackage{graphicx}
\usepackage{bm}              % bold math
\usepackage{xcolor}
\usepackage{hyperref}
\usepackage{siunitx}         % proper SI units
\usepackage{booktabs}        % nice tables

% ── convenience macros ────────────────────────────────────────────────────────
\newcommand{\dd}{\mathrm{d}}               % upright differential
\newcommand{\vA}{v_{\!A}}                  % self-propulsion speed
\newcommand{\peff}{p_{\mathrm{eff}}}       % effective shape index
\newcommand{\Jk}{\tilde{J}}                % dimensionless adhesion J/(2γ)
\newcommand{\phii}{\phi_i}                 % cell field shorthand
\newcommand{\phij}{\phi_j}                 % neighbor field shorthand
\newcommand{\todo}[1]{{\color{red}\textbf{[TODO: #1]}}}

% ══════════════════════════════════════════════════════════════════════════════
\begin{document}

\title{Adhesion-controlled rigidity transition in a multi-cell phase field model}

\author{Steven~A.~Silber}
\affiliation{Department of Physics and Astronomy, Western University, 1151 Richmond Street, London, Ontario, Canada N6A 3K7}

\author{Mikko~Karttunen}
\email{mikko.karttunen1@uef.fi}
\affiliation{Department of Physics and Astronomy, Western University, 1151 Richmond Street, London, Ontario, Canada N6A 3K7}
\affiliation{ELLIS Institute Finland, Maarintie 8, 02150 Espoo, Finland}
\affiliation{Department of Technical Physics, University of Eastern Finland, P.O.\ Box 1627, FI-70211 Kuopio, Finland}
\affiliation{Department of Chemistry, Western University, 1151 Richmond Street, London, Ontario, Canada N6A 5B7}

\date{\today}

% ── abstract ──────────────────────────────────────────────────────────────────
\begin{abstract}
The vertex model predicts a rigidity transition at a critical target shape index
$p_0^*\!\approx 3.81$, controlled by the balance between cortical contractility
and cell-cell adhesion.  Whether an analogous transition exists in continuum
phase field descriptions---where cell boundaries are diffuse and topology is
unconstrained---remains an open question.  Here we incorporate gradient-coupling
adhesion,
$F_{\mathrm{adh}} = J\sum_{i<j}\!\int\!\nabla\phii\cdot\nabla\phij\,\dd A$,
into a GPU-accelerated, multi-cell phase field model of a confluent monolayer
($N = 288$--$1152$ cells, packing fraction $\phi = 0.89$) and derive the
stability bound~$J < 2\gamma$.  The dimensionless parameter
$\Jk = J/(2\gamma)$ measures the fraction of surface energy removed at shared
interfaces.  Adhesion quench experiments---in which adhesion is instantaneously
enabled from an equilibrated zero-adhesion state---reveal a sharp rigidity
transition: below~$\Jk \approx 0.25$ cells remain caged, while
above~$\Jk \approx 0.50$ they undergo full rearrangements driven purely by
overdamped relaxation without self-propulsion.  This demonstrates that
gradient-coupling adhesion tunes the energy landscape in a manner analogous to
the shape index in vertex models, establishing a path toward quantitative
comparison between the two frameworks.
\end{abstract}

\maketitle

% ══════════════════════════════════════════════════════════════════════════════
%  I.  INTRODUCTION
% ══════════════════════════════════════════════════════════════════════════════
\section{Introduction}
\label{sec:intro}

Epithelial monolayers can switch between a solid-like, mechanically rigid
state and a fluid-like, collectively migrating one.  This jamming transition
has been observed in wound-edge MDCK
monolayers~\cite{Angelini2011,Garcia2015}, in bronchial epithelial cells
from asthmatic donors~\cite{Park2015}, along the vertebrate body
axis~\cite{Mongera2018}, and in a variety of other epithelial
systems~\cite{Atia2018,Malinverno2017}.  The experimental evidence
consistently points to cell shape as the relevant order parameter: unjammed
monolayers contain elongated cells with a shape index $p = P/\sqrt{A}$
exceeding $\approx 3.81$, while jammed tissues consist of compact
cells~\cite{Park2015,Atia2018}.

The theoretical picture has been built largely on the vertex model, in
which each cell is a polygon with energy
$E = \sum_i[K_A(A_i - A_0)^2 + K_P(P_i - P_0)^2]$, and the target shape
index $p_0 = P_0/\sqrt{A_0}$ serves as the control
parameter~\cite{NagaiHonda2001,Farhadifar2007,Staple2010,Fletcher2014}.
Bi \textit{et al.}~\cite{Bi2015} showed that a rigidity transition occurs at
$p_0^* \approx 3.81$, and with added self-propulsion the transition extends
to a full $(p_0,\,v_0)$ phase diagram~\cite{Bi2016}.  Subsequent work
introduced active polarity~\cite{Barton2017}, multicellular
rosettes~\cite{YanBi2019}, 3D shape criteria~\cite{MerkelManning2018},
and cell turnover~\cite{Czajkowski2019}.  In the vertex model, adhesion
enters indirectly: cadherin-mediated adhesion reduces the line tension at
cell-cell junctions, which raises $p_0$ and pushes the tissue toward the
fluid phase.

Phase field models offer a complementary description in which each cell~$i$
is represented by a continuous scalar field $\phii(\bm{r},t)$, unity inside
the cell and zero outside, with a diffuse interface of width~$\lambda$.
Rearrangements occur through smooth field evolution rather than discrete
topology changes, and cell shapes are unconstrained by a polygonal
representation~\cite{Camley2017,AlertTrepat2020}.  Single-cell phase field
models were introduced by Shao \textit{et al.}~\cite{Shao2010} and Ziebert
\textit{et al.}~\cite{Ziebert2012}; Nonomura~\cite{Nonomura2012} developed
the first multi-cell formulation.

The treatment of cell-cell adhesion in multi-cell phase field models varies
considerably.
Nonomura~\cite{Nonomura2012} introduced a gradient-overlap coupling
$-\gamma\sum_{i<j}\!\int\! \nabla h(\phii)\!\cdot\!\nabla h(\phij)\,\dd A$,
which is localized at the diffuse interface and acts as an effective cortical
adhesion.
L\"ober \textit{et al.}~\cite{Lober2015} used quartic overlap
$\lambda\sum_{i<j}\!\int\!\rho_i^2\rho_j^2\,\dd A$ for steric repulsion, and modeled adhesion as a
non-variational advection that drives each cell's interface along the
outward normal of its neighbor.
Najem and Grant~\cite{Najem2016} coupled the cell cortex to an auxiliary
range field~$C_i$, producing an adhesion with independently controllable
range; they used this to measure tissue surface tension.
Palmieri \textit{et al.}~\cite{Palmieri2015} employed quartic repulsion
only.
More recent studies have likewise omitted an explicit adhesion term and
focused on other questions:
Loewe \textit{et al.}~\cite{Loewe2020} studied the solid-liquid transition
as a function of cell deformability and motility,
Wenzel and Voigt~\cite{Wenzel2021} compared four activity models using a
signed-distance-based repulsion,
Graham \textit{et al.}~\cite{Graham2024} demonstrated activity-driven cell
sorting without adhesion, and
Saito and Ishihara~\cite{Saito2024} showed that cell deformability alone
can drive a fluid-to-fluid transition.
To date, no multi-cell phase field study has systematically varied adhesion
strength across a rigidity transition.

In this work, we introduce gradient-coupling adhesion
%
\begin{equation}
  F_{\mathrm{adh}} = J\sum_{i<j}\int \nabla\phii\cdot\nabla\phij\;\dd A,
  \label{eq:adhesion-intro}
\end{equation}
%
into a multi-cell phase field model.  At a shared interface where two
cells meet, their field gradients are anti-parallel
($\nabla\phii \approx -\nabla\phij$), so the integrand is negative and
the energy is lowered.  Away from interfaces, both gradients vanish and
the term contributes nothing: the adhesion is strictly
surface-localized, encoding the same contact-length-dependent physics
as cadherin-mediated adhesion in vertex models.  The dimensionless
control parameter $\Jk = J/(2\gamma)$ measures the fraction of surface
energy removed at shared interfaces, establishing the natural analog
of the vertex model target shape index~$p_0$.

We probe the transition using three protocols.
(i)~An \emph{adhesion quench} (Sec.~\ref{sec:quench}) at zero motility,
in which adhesion is switched on instantaneously from an equilibrated
repulsive state.  This protocol has no vertex model analog, since at
$T = 0$ the vertex model is frozen regardless of
$p_0$~\cite{Bi2015}; the phase field, by contrast, relaxes via gradient
descent, exposing the continuous trajectory of rearrangement events.
(ii)~A \emph{motility probe} (Sec.~\ref{sec:motility}), in which we sweep
$\Jk$ at small fixed $\vA$ to locate the jammed-to-fluid boundary.
(iii)~A full \emph{$(\Jk,\,\vA)$ phase diagram}
(Sec.~\ref{sec:diagram}), compared to the $(p_0, v_0)$ diagram of
Bi \textit{et al.}~\cite{Bi2016}.

The model and equations of motion are defined in Sec.~\ref{sec:model},
simulation methods in Sec.~\ref{sec:methods}, results in
Secs.~\ref{sec:quench}--\ref{sec:diagram}, and discussion and conclusions
in Secs.~\ref{sec:discussion}--\ref{sec:conclusion}.


% ══════════════════════════════════════════════════════════════════════════════
%  II.  MODEL
% ══════════════════════════════════════════════════════════════════════════════
\section{Model}
\label{sec:model}

% --- II.A  Multi-cell phase field formulation ---
\subsection{Multi-cell phase field}
\label{sec:model:pf}

We consider $N$ cells on a two-dimensional square domain of side~$L$ with
periodic boundary conditions.  Each cell $i = 1,\dots,N$ is described by a
scalar order parameter $\phii(\bm{r},t) \in [0,1]$, where $\phii \approx 1$
inside the cell and $\phii \approx 0$ outside.  The total free energy is
%
\begin{equation}
  F[\{\phii\}] = \sum_i F_i^{\mathrm{single}}
                 + F^{\mathrm{rep}}
                 + F^{\mathrm{adh}},
  \label{eq:total-energy}
\end{equation}
%
where the three contributions are defined below.

\paragraph{Single-cell energy.}
Each cell has a gradient (interfacial) and a bulk (double-well) contribution:
%
\begin{equation}
  F_i^{\mathrm{single}}
    = \int \!\biggl[
        \frac{\epsilon^2}{2}\,|\nabla\phii|^2
        + \frac{1}{4}\,\phii^2\,(1 - \phii)^2
      \biggr]\dd A
    + \mu\biggl(\int\!\phii\;\dd A - V_0\biggr)^{\!2}.
  \label{eq:single-cell}
\end{equation}
%
The gradient term penalizes sharp interfaces with a stiffness set by the
interface width parameter~$\epsilon$ (related to~$\lambda$, the half-width of
the $\tanh$ profile).  The double-well potential has minima at $\phii = 0$
(exterior) and $\phii = 1$ (interior), stabilizing a sharp-interface profile
of the form $\phii \approx \tfrac{1}{2}[1 - \tanh(d/\lambda)]$, where $d$ is
the signed distance from the cell boundary.  The quadratic volume constraint
with Lagrange-like penalty~$\mu$ maintains the cell area near a target
value~$V_0 = \pi R^2$, where $R$ is the nominal cell radius.

\paragraph{Repulsive interaction.}
Overlap between distinct cells is penalized by a quartic term:
%
\begin{equation}
  F^{\mathrm{rep}} = \kappa \sum_{i<j} \int \phii^2\,\phij^2\;\dd A.
  \label{eq:repulsion}
\end{equation}
%
This steric repulsion, introduced by Nonomura~\cite{Nonomura2012} in the
first multi-cell phase field model, prevents cell merger and sets the energy
scale for intercellular interactions.  Because $\phii^2\phij^2$ grows as the
fourth power of the fields, the cost of overlap rises steeply once two cells'
interiors begin to coincide.

\paragraph{Adhesive interaction.}
Three distinct adhesion mechanisms have appeared in the multi-cell phase
field literature, each shaped by the physical question its authors sought
to answer.

Nonomura~\cite{Nonomura2012} designed the first multi-cell phase field
model around the smooth step function $h(\phi) = \phi^2(3 - 2\phi)$, whose
gradient $\nabla h(\phii)$ is localized at the diffuse interface and
serves as a mathematical representation of the cell cortex.
Adhesion is a gradient-overlap coupling
$-\gamma\sum_{i<j}\!\int\!\nabla h(\phii)\!\cdot\!\nabla h(\phij)\,\dd A$:
it is nonzero only where two cortices overlap, directly encoding
cadherin-mediated surface adhesion.  Nonomura states that this form was
chosen because it is ``the simplest among the several alternatives'' that
can be rewritten in terms of the per-cell-type sum field
$w_\ell = \sum_m h(\phi_m)\,\delta_{\ell_m,\ell}$, enabling efficient
$O(NL)$ evaluation~\cite{Nonomura2012}.  The model was applied
to cell sorting---a differential-adhesion problem---not to a rigidity
transition.

L\"ober, Ziebert, and Aranson~\cite{Lober2015} extended their single-cell
actin-substrate motility model~\cite{Ziebert2012} to many cells.  Because
the underlying model already contains non-variational terms---actin
advection, substrate friction, motor contractility---interactions are added
at the equation-of-motion level rather than through a free energy.  Their
adhesion term
$-\kappa\,\nabla\rho_i \!\cdot\! \hat{f}(\nabla\rho_j)$
advects cell~$i$'s boundary along the outward normal of cell~$j$, mimicking
the mechanical attraction of adjacent cell membranes~\cite{Lober2015}.
This formulation cannot, by construction, be derived from a free energy
functional.

Najem and Grant~\cite{Najem2016} began from an explicit continuum
mechanics picture: an exponentially decaying adhesion potential
$W(x) = w\exp(-d^2/\delta^2)$ between cell surfaces.  They translated
this into phase field language via an auxiliary range field~$C_i$ that
defines each cell's interaction neighborhood independently of the
interface width.  The resulting adhesion energy
$w\,C_i\,\phii^2(1 - \phii)\sum_{j\neq i}(1 - \phij)$
acts at cell~$i$'s cortex ($\phii^2(1-\phii)$ is localized at the
interface), reaches into the intercellular gap via~$(1 - \phij)$, and has a
range set by~$C_i$ rather than by~$\epsilon$.  This decoupling of range
from interface width was essential for their goal: measuring tissue surface
tension as a function of adhesion-to-cortical-tension ratio,
$T_{\mathrm{st}}/\sigma = 2\,w/\sigma$.

The remaining multi-cell phase field studies---Palmieri
\textit{et al.}~\cite{Palmieri2015},
Loewe \textit{et al.}~\cite{Loewe2020},
Wenzel and Voigt~\cite{Wenzel2021},
Graham \textit{et al.}~\cite{Graham2024}---omit adhesion entirely,
using only quartic repulsion while varying motility, deformability,
or activity as control parameters.  To date, no multi-cell phase field
study has systematically varied adhesion strength across a rigidity
transition.

A common feature of the three adhesion models is that each localizes the
interaction at the cell interface: the gradient in Nonomura, the boundary
normal in L\"ober, and the cortex factor
$\phii^2(1 - \phii)$ in Najem--Grant.  This reflects the biological
picture that cadherins are transmembrane proteins on the cell surface.
However, the standard quartic \emph{repulsion}
$\kappa\,\phii^2\phij^2$ is itself a bulk overlap term---it is
nonzero wherever both fields coexist, not only at the interface---yet
it is universally adopted without objection.  The expectation that
adhesion must be interface-localized while repulsion need not be
is a modeling convention, not a physical requirement: in the
sharp-interface limit ($\epsilon/R \to 0$), both bulk and
interface-localized couplings reduce to energies proportional to the
pairwise contact length~$\ell_{ij}$.

\paragraph{Choice of the bilinear form.}
To fill this gap with minimal assumptions, we expand the most general
pairwise coupling to lowest orders:
%
\begin{equation}
  F^{\mathrm{pair}}_{ij}
    = \int\!\bigl[
        a_{11}\,\phii\,\phij
        + a_{22}\,\phii^2\,\phij^2
        + \cdots
      \bigr]\dd A,
  \label{eq:landau-pair}
\end{equation}
%
where cubic terms ($\phii^2\phij$) are excluded by $i \leftrightarrow j$
exchange symmetry.  The quartic coefficient is fixed by the repulsion
($a_{22} = \kappa > 0$); the bilinear coefficient~$a_{11}$ is the only
remaining free parameter at leading order.  Setting $a_{11} = -J$ with
$J \geq 0$ yields an attraction that rewards overlap between neighboring
fields---the natural lowest-order complement to the quartic repulsion,
analogous to the relationship between a harmonic attraction and a
hard-core repulsion in molecular dynamics.  We therefore adopt
%
\begin{equation}
  F^{\mathrm{adh}} = -J\sum_{i<j}\int \phii\,\phij\;\dd A,
  \label{eq:adhesion}
\end{equation}
%
where $J \geq 0$ is the adhesion coupling strength.  The variational
derivative with respect to cell~$i$ is
%
\begin{equation}
  \frac{\delta F^{\mathrm{adh}}}{\delta \phii}
    = -J \sum_{j \neq i} \phij,
  \label{eq:adhesion-var}
\end{equation}
%
which is simply the total phase field of all other cells evaluated at
position~$\bm{r}$---the simplest possible intercellular coupling.

Unlike the gradient-based forms of Nonomura~\cite{Nonomura2012} and
L\"ober~\cite{Lober2015}, which are strictly interface-localized, the
bilinear coupling is nonzero wherever both fields are nonzero and
thus extends into the cell body.
In the sharp-interface limit, however, both forms converge to an
energy proportional to the pairwise contact length.  For two cells
with $\tanh$-profile interfaces of width~$\epsilon$, the overlap
integral scales as
%
\begin{equation}
  \int\!\phii\,\phij\;\dd A
    \;\sim\;\epsilon\,\ell_{ij},
  \label{eq:sharp-interface}
\end{equation}
%
where $\ell_{ij}$ is the contact length between cells~$i$ and~$j$.
The bilinear adhesion energy is therefore
$F^{\mathrm{adh}} \sim -J\epsilon\sum_{i<j}\ell_{ij}$,
which is the continuum analog of the vertex model adhesion energy
$-\gamma\sum_{i<j}\ell_{ij}$ with the identification
$\gamma = J\epsilon$.  This sharp-interface correspondence
establishes that the bilinear coupling is a legitimate encoding of the
same contact-length-dependent adhesion that appears in vertex models.

\paragraph{Pair-energy analysis and the control parameter $\Jk$.}
To see how $\Jk$ controls cell-cell overlap, consider two cells at
a point of symmetric overlap $\phii = \phij \equiv \phi$.  The local
pairwise energy density from Eqs.~\eqref{eq:repulsion}
and~\eqref{eq:adhesion} is
%
\begin{equation}
  e(\phi) = \kappa\,\phi^4 - J\,\phi^2,
  \label{eq:pair-energy}
\end{equation}
%
which has a minimum at $\phi_{\mathrm{eq}}^2 = J/(2\kappa)$ with depth
$e_{\min} = -J^2/(4\kappa)$.  At $J = 0$, the minimum is at~$\phi = 0$
(no overlap---hard repulsion), so the purely repulsive model
\cite{Palmieri2015,Loewe2020} puts cell interfaces as far apart as density
allows.  As $\Jk$ increases, $\phi_{\mathrm{eq}}$ rises: cells settle into
a progressively deeper overlap well, widening their shared interface and
increasing their contact area.  Because larger contact area between
neighbors requires more elongated, higher-perimeter cells, increasing~$\Jk$
drives cells toward larger effective shape indices---exactly the mechanism by
which rising~$p_0$ drives the vertex model toward the fluid phase.
We therefore treat the single dimensionless ratio
%
\begin{equation}
  \Jk
  \label{eq:jk-ratio}
\end{equation}
%
as the primary control parameter throughout this work, and argue that it
plays the role of the vertex model target shape index.

\paragraph{Stability.}
Cell merger (complete overlap of two fields) requires
$\kappa\,\phii^2\phij^2 < J\,\phii\,\phij$ throughout the overlap region.
In the cell interior where $\phii \approx \phij \approx 1$, this reduces to
$\kappa < J$, i.e., $\Jk > 1$, setting a hard upper
bound~\cite{Nonomura2012}.  At the diffuse interface where
$\phii,\,\phij < 1$, the bilinear adhesion dominates the quartic repulsion
at lower~$J$, so the \emph{effective} stability boundary is
$\Jk \lesssim 0.5$--$1$, depending on the interface profile and packing
fraction.

Prior phase field studies have explored adhesion values only qualitatively:
Nonomura~\cite{Nonomura2012} tested $\gamma/\kappa$ values of order $10^{-2}$;
Najem and Grant~\cite{Najem2016}, using their range-field adhesion,
varied the ratio of adhesion to cortical tension $w/\sigma$ over~$[0.1,\,0.5]$
and found the tissue surface tension scales linearly,
$T_{\mathrm{st}}/\sigma = 2\,w/\sigma$, confirming the differential adhesion
hypothesis in two dimensions.
In this work, we scan a range $\Jk \in [0,\,0.5]$ chosen to approach
but not exceed the stability limit, and we determine the empirical merger
boundary in Sec.~\ref{sec:quench}.


% --- II.B  Equations of motion ---
\subsection{Equations of motion}
\label{sec:model:eom}

The phase field evolves via overdamped (model-A) dynamics with an additional
self-propulsion term:
%
\begin{equation}
  \frac{\partial \phii}{\partial t}
    = -\frac{\delta F}{\delta \phii}
      + \alpha\,\vA^{(i)}\,\hat{\bm{p}}_i \cdot \nabla\phii,
  \label{eq:eom}
\end{equation}
%
where $\alpha$ is a motility coupling constant, $\vA^{(i)}$ is the
self-propulsion speed of cell~$i$, and $\hat{\bm{p}}_i$ is its polarity unit
vector.  The polarity undergoes rotational diffusion (run-and-tumble
dynamics):
%
\begin{equation}
  \frac{\dd\hat{\bm{p}}_i}{\dd t}
    = \frac{1}{\tau}\bigl(\hat{\bm{n}}_i - \hat{\bm{p}}_i\bigr),
  \label{eq:polarity}
\end{equation}
%
where $\tau$ is the persistence (tumble) time and $\hat{\bm{n}}_i$ is a
random unit vector that is resampled at Poisson-distributed times with
rate~$1/\tau$.  In the zero-motility limit ($\vA = 0$) used in the quench
protocol, the dynamics reduce to pure gradient descent on the free
energy~$F$.

The cell centroid velocity is computed from the first moment of the field
evolution:
%
\begin{equation}
  \bm{v}_i
    = \frac{\displaystyle\int (\partial_t \phii)\,\bm{r}\;\dd A}
           {\displaystyle\int \phii\;\dd A}.
  \label{eq:velocity}
\end{equation}

\paragraph{Connection to the vertex model.}
In the vertex model, the target shape index $p_0 = P_0/\sqrt{A_0}$ is a
\emph{control parameter} in the Hamiltonian: it sets
the preferred perimeter-to-area ratio, and the rigidity transition occurs
at $p_0^* \approx 3.81$.  The actual cell shape index
$p = P/\sqrt{A}$ is an \emph{emergent observable} that generally differs
from~$p_0$ and must be measured from the resulting configurations.
Adhesion at cell--cell interfaces lowers the effective line tension and
raises the target~$p_0$.

In our phase field model, the ratio~$\Jk$ is the control parameter: it
appears in the free energy~\eqref{eq:total-energy} and is set before the
simulation.  The adhesion integral $\int\!\phii\,\phij\;\dd A$ measures the
contact area between cells~$i$ and~$j$, which is the continuum analog of
the shared edge length~$\ell_{ij}$ in the vertex model.  Increasing~$J$
energetically favors larger contact area, which deforms cells toward more
elongated, larger-perimeter shapes---precisely the effect of raising~$p_0$.
The effective shape index~$\peff$~\eqref{eq:peff}, extracted from the
$\phii = 0.5$ contour, is the emergent observable that we \emph{measure}
from the resulting configurations.  The central question of this work is
whether the rigidity transition in~$\Jk$ coincides with
$\langle\peff\rangle \approx 3.81$---i.e., whether the vertex model's
geometric criterion for the transition extends to the continuum phase field
framework.  The mapping $\Jk \to \peff$ is therefore a \emph{result} of
the simulation, not an assumption.


% ══════════════════════════════════════════════════════════════════════════════
%  III.  SIMULATION METHODS
% ══════════════════════════════════════════════════════════════════════════════
\section{Simulation methods}
\label{sec:methods}

% --- III.A  Implementation ---
\subsection{GPU-accelerated implementation}
\label{sec:methods:gpu}

The model is implemented as a custom CUDA C++ code, building on the
GPU-accelerated phase field framework described in
Ref.~\cite{Silber2025}.  All $N$ phase fields are
stored on a single GPU and evolved simultaneously.  The per-cell fields
$\phii$ are represented on a uniform Cartesian grid of $L \times L$ points
with grid spacing $\Delta x = 1$.  At each time step a \emph{scatter} kernel
accumulates the auxiliary sum fields
%
\begin{equation}
  S_2(\bm{r}) = \textstyle\sum_k \phi_k^2(\bm{r}), \qquad
  S_1(\bm{r}) = \textstyle\sum_k \phi_k(\bm{r})
  \label{eq:sum-fields}
\end{equation}
%
over all cells, and then a \emph{fused} kernel evaluates the variational
derivative, interaction terms, self-propulsion, and forward-Euler update for
every cell and grid point in a single pass.  The repulsive force on cell~$i$
at position~$\bm{r}$ is $2\kappa\,\phii\bigl(S_2 - \phii^2\bigr)$; the
adhesive force is $-J\bigl(S_1 - \phii\bigr)$.  When $J = 0$, the adhesion
field~$S_1$ is neither allocated nor computed, so the code incurs zero overhead
for the purely repulsive case.

The implementation achieves \todo{X}~time units per second for $N = 288$ cells
on an NVIDIA H100 MIG partition (10~GB), and \todo{X}~t/s on an RTX~4090
laptop GPU.  All production runs were performed on the Nibi cluster
(Alliance Canada).

% --- III.B  Parameters ---
\subsection{Simulation parameters}
\label{sec:methods:params}

Table~\ref{tab:params} summarizes the physical and numerical parameters, which
are held fixed across all adhesion experiments.

\begin{table}[b]
  \caption{%
    Fixed simulation parameters.  The time scale is set by $\tau = 10{,}000$
    (tumble time); the length scale by the cell radius $R = 49$.  The packing
    fraction $\phi = N \pi R^2 / L^2 = 0.89$ places the system in the
    confluent regime.
  }
  \label{tab:params}
  \begin{ruledtabular}
  \begin{tabular}{lll}
    Parameter & Symbol & Value \\ \hline
    Cell count       & $N$        & 288 (or 1152)     \\
    Domain side      & $L$        & 1562 (or 3124)    \\
    Cell radius      & $R$        & 49                \\
    Packing fraction & $\phi$     & 0.89              \\
    Interface width  & $\lambda$  & 7                 \\
    Repulsion        & $\kappa$   & 10                \\
    Time step        & $\Delta t$ & 0.02              \\
    Tumble time      & $\tau$     & $10{,}000$        \\
    Boundary cond.   &            & periodic          \\
  \end{tabular}
  \end{ruledtabular}
\end{table}

The packing fraction $\phi = 0.89$ was chosen to match the near-confluent
regime in which the vertex model is defined.  At this density, cells tile
the domain with only small interstitial gaps, consistent with a confluent
monolayer.

% --- III.C  Equilibration ---
\subsection{Equilibration protocol}
\label{sec:methods:eq}

All adhesion experiments start from a well-equilibrated, purely repulsive
configuration.  One hundred independent realizations of 288 cells were
initialized at random positions on the $L = 1562$ domain and relaxed at
$\vA = 0$, $J = 0$ for $t = 80{,}000$ ($8\tau$), following the equilibration
protocol of Palmieri \textit{et al.}~\cite{Palmieri2015}.  At this point, all
transient stress from the random initialization has decayed and the system
sits in a local energy minimum of the repulsive free energy.  The resulting
100 checkpoints serve as independent starting states for the adhesion
experiments.

To verify equilibration quality, a control run at $J = 0$, $\vA = 0$ is
included in every experiment: a well-equilibrated starting state should
produce negligible centroid displacement over the measurement window.

% --- III.D  Observables ---
\subsection{Observables}
\label{sec:methods:obs}

\paragraph{Mean-squared displacement (MSD).}
The ensemble- and time-averaged MSD is defined as
%
\begin{equation}
  \langle \Delta r^2(t) \rangle
    = \biggl\langle \frac{1}{N}\sum_{i=1}^{N}
      |\bm{r}_i(t_0 + t) - \bm{r}_i(t_0)|^2 \biggr\rangle_{t_0},
  \label{eq:msd}
\end{equation}
%
where the angular brackets denote an average over starting times~$t_0$.
Displacements are unwrapped across periodic boundaries.  In the jammed phase,
$\langle \Delta r^2 \rangle$ plateaus at a value set by the cage size; in the
fluid phase it grows linearly, $\langle \Delta r^2 \rangle \sim 4 D_{\mathrm{eff}}\,t$.

\paragraph{Non-Gaussian parameter.}
The non-Gaussian parameter
%
\begin{equation}
  \alpha_2(t)
    = \frac{\langle \Delta r^4(t)\rangle}
           {2\,\langle \Delta r^2(t)\rangle^2} - 1
  \label{eq:alpha2}
\end{equation}
%
measures dynamic heterogeneity.  A peak in $\alpha_2$ at time $t^*$ signals
a population of fast-moving cells coexisting with caged cells, characteristic
of the approach to the glass/jamming transition.

\paragraph{Cage-relative MSD.}
To isolate local rearrangements from collective drift, we compute the
cage-relative displacement~\cite{Bi2016}:
%
\begin{equation}
  \Delta \bm{r}_i^{\,\mathrm{cage}}(t)
    = \Delta\bm{r}_i(t)
      - \frac{1}{n_i}\sum_{j \in \mathrm{nn}(i)} \Delta\bm{r}_j(t),
\end{equation}
%
where the sum runs over the $n_i$ nearest neighbors of cell~$i$ at $t = 0$.

\paragraph{Effective shape index.}
From VTK snapshots of the phase fields, we extract the $\phii = 0.5$ contour
for each cell using a marching-squares algorithm, then compute the perimeter
$P_i$ and area $A_i$ of the contour.  The effective shape index is
%
\begin{equation}
  \peff = \frac{P_i}{\sqrt{A_i}},
  \label{eq:peff}
\end{equation}
%
which can be compared directly to the vertex model quantity~$p_0$.

\paragraph{Contact area.}
The pairwise contact area between cells $i$ and $j$ is
%
\begin{equation}
  C_{ij} = \int \phii\,\phij\;\dd A,
  \label{eq:contact-area}
\end{equation}
%
which is the integrand of the adhesion energy~\eqref{eq:adhesion}.  The total
contact area $C_{\mathrm{tot}} = \sum_{i<j} C_{ij}$ measures the degree of
cell--cell contact and is expected to increase with~$\Jk$.

\paragraph{Energy decomposition.}
The free energy~\eqref{eq:total-energy} is decomposed into gradient, bulk,
repulsive, and adhesive contributions, each tracked as a function of time
during every simulation.  In the quench protocol (Sec.~\ref{sec:quench}),
the time series of each component distinguishes interface relaxation (smooth
gradient energy decay) from topological rearrangements (stepwise drops in
interaction energy).

\paragraph{Neighbor topology.}
Two cells $i$ and $j$ are defined as neighbors if $C_{ij}$ exceeds a threshold
$C_{\min}$.  Changes in the neighbor list between the start and end of a
simulation quantify the number of T1-like rearrangement events.

% --- III.E  Protocol summary ---
\subsection{Experimental protocols}
\label{sec:methods:protocols}

Three experimental phases, summarized in Table~\ref{tab:protocols}, probe the
rigidity transition at increasing levels of detail.

\begin{table}[t]
  \caption{%
    Summary of the three experimental protocols.  All start from the same
    pool of 100 independently equilibrated checkpoints at $t = 80{,}000$.
  }
  \label{tab:protocols}
  \begin{ruledtabular}
  \begin{tabular}{lccc}
    & Phase 0     & Phase 1       & Phase 2           \\
    & (quench)    & (probe)       & (diagram)         \\ \hline
    $\Jk$
      & 0--0.5 (8 pts)  & 0--0.5 (8 pts) & 0--0.5 (6 pts) \\
    $\vA$
      & 0               & 0.002          & 0.002--0.012    \\
    Duration
      & $2\tau$          & $5\tau$        & $5\tau$         \\
    Replicates
      & 1 (deterministic)& 3             & 3--10           \\
    Key observable
      & displacement,    & MSD,           & MSD,            \\
      & energy           & $\alpha_2$     & $\peff$         \\
  \end{tabular}
  \end{ruledtabular}
\end{table}


% ══════════════════════════════════════════════════════════════════════════════
%  IV.  RESULTS — Phase 0: Adhesion quench
% ══════════════════════════════════════════════════════════════════════════════
\section{Adhesion quench at zero motility}
\label{sec:quench}

% Narrative plan:
% 1. Explain protocol: instantaneous J-quench from equilibrated J=0 state
% 2. Show total centroid displacement vs J/κ — the key figure
%    - Expect: flat at low J/κ (interface adjustments only), rising above (J/κ)*
% 3. Energy time series at representative J/κ values
%    - Below (J/κ)*: single smooth decay (gradient energy relaxation)
%    - Above (J/κ)*: staircase / multi-step decay (rearrangement events)
% 4. Neighbor topology changes: count T1-like events vs J/κ
% 5. Relaxation timescale vs J/κ — should diverge at (J/κ)*
% 6. Determine empirical (J/κ)* and stability limit

\todo{Awaiting Phase 0 data from cluster (jobs 8661054--8661061).  The
following subsections outline the planned analysis and expected results.}

\subsection{Protocol}
\label{sec:quench:protocol}

Starting from a checkpoint equilibrated at $J = 0$, $\vA = 0$ for
$t = 80{,}000$ ($8\tau$), we instantaneously set $J > 0$ while keeping
$\vA = 0$ and evolve the system for an additional $2\tau$.  The dynamics are
purely relaxational: $\partial_t \phii = -\delta F/\delta\phii$.  Because
the $J = 0$ equilibrium is generally \emph{not} a minimum of the
$J > 0$ energy landscape, the system flows downhill to a new minimum.

This ``$J$-quench'' has no vertex model analog.  In the vertex model at
$T = 0$, $v_0 = 0$, the system is frozen in a local minimum regardless
of~$p_0$; changing $p_0$ simply relabels the energy of existing configurations
without inducing motion.  In the phase field, turning on adhesion creates an
energetic driving force and the overdamped dynamics trace out the
\emph{continuous} relaxation path---including intermediate configurations
that would be inaccessible in a discrete-topology model.

\subsection{Centroid displacement}
\label{sec:quench:displacement}

\todo{Figure: total displacement $\sum_i |\Delta\bm{r}_i|$ vs.\ $\Jk$.
Expected: small displacement (interface adjustment) for $\Jk < (\Jk)^*$,
large displacement (cell rearrangements) for $\Jk > (\Jk)^*$.  The control
($\Jk = 0$) should show near-zero displacement, confirming equilibration
quality.}

\subsection{Energy landscape}
\label{sec:quench:energy}

\todo{Figure: energy components (gradient, bulk, repulsion, adhesion) vs.\
time for representative $\Jk$ values.  Expected: smooth monotonic decay
below $(\Jk)^*$; stepwise or multi-plateau decay above $(\Jk)^*$,
corresponding to discrete rearrangement events resolved continuously
in time.}

\subsection{Neighbor rearrangements}
\label{sec:quench:topology}

\todo{Count T1-like neighbor changes (new contacts formed, old contacts lost)
between start and end of each quench.  Plot vs.\ $\Jk$.}

\subsection{Relaxation timescale}
\label{sec:quench:timescale}

\todo{Define relaxation time from energy decay or displacement saturation.
Plot vs.\ $\Jk$; expect divergence near $(\Jk)^*$ (critical slowing down).}


% ══════════════════════════════════════════════════════════════════════════════
%  V.  RESULTS — Phase 1: Motility probe
% ══════════════════════════════════════════════════════════════════════════════
\section{Motility probe sweep}
\label{sec:motility}

% Narrative plan:
% 1. At fixed small v_A, sweep J/κ
% 2. MSD: plateau (jammed) vs linear (fluid)
% 3. α₂: peak height and position vs J/κ
% 4. Cage-relative MSD
% 5. Classify jammed/fluid → locate (J/κ)* at this v_A
% 6. Compare to Phase 0 result

\todo{Phase 1 runs not yet submitted.  Below is the analysis plan.}

\subsection{Mean-squared displacement}
\label{sec:motility:msd}

\todo{Figure: MSD vs.\ time for each $\Jk$ at $\vA = 0.002$.  Low $\Jk$
should show caging (plateau); high $\Jk$ should show diffusive scaling.
Transition point identifies $(\Jk)^*$ at this motility.}

\subsection{Dynamic heterogeneity}
\label{sec:motility:alpha2}

\todo{Figure: $\alpha_2(t)$ for each $\Jk$.  Peak at $t^* \sim 2\tau$
expected near the transition, suppressed deep in jammed or fluid regimes.}

\subsection{Effective shape index}
\label{sec:motility:shape}

\todo{Extract $\peff$ from VTK frames.  Plot $\langle\peff\rangle$ vs.\
$\Jk$; compare to the vertex model critical value $p_0^* \approx 3.81$.
This directly tests whether $\Jk$ maps onto the shape index.}


% ══════════════════════════════════════════════════════════════════════════════
%  VI.  RESULTS — Phase 2: Phase diagram
% ══════════════════════════════════════════════════════════════════════════════
\section{Adhesion--motility phase diagram}
\label{sec:diagram}

% Narrative plan:
% 1. Heatmap: MSD at t = 5τ in (J/κ, v_A) plane
% 2. Phase boundary from MSD threshold (or cage-relative crossover)
% 3. Comparison to Bi (2016) (p_0, v_0) diagram
% 4. Shape index at each (J/κ, v_A) point → overlay on diagram
% 5. System-size comparison: 288 vs 1152 cells

\todo{Phase 2 runs not yet designed in detail.  Shell for the phase diagram
analysis.}

\subsection{Phase boundary}
\label{sec:diagram:boundary}

\todo{Figure: heatmap of $\langle\Delta r^2(5\tau)\rangle$ in the
$(\Jk,\,\vA)$ plane.  Draw phase boundary separating jammed (MSD plateau)
from fluid (MSD growing).  Compare shape to the vertex model
$(p_0,\,v_0)$ boundary of Bi \textit{et al.}~\cite{Bi2016}.}

\subsection{Shape index mapping}
\label{sec:diagram:shape}

\todo{At each $(\Jk,\,\vA)$ point, compute $\langle\peff\rangle$.  Overlay
contours of constant $\peff$ on the phase diagram.  If the phase boundary
coincides with $\peff \approx 3.81$, the vertex model mapping is quantitative.}

\subsection{System-size effects}
\label{sec:diagram:size}

\todo{Compare 288-cell and 1152-cell results at selected points.  Check
whether finite-size effects shift the transition.}


% ══════════════════════════════════════════════════════════════════════════════
%  VII.  DISCUSSION
% ══════════════════════════════════════════════════════════════════════════════
\section{Discussion}
\label{sec:discussion}

% Key discussion points:

% 1. J/κ as the phase field analog of p_0
%    - Both control the balance between cell compactness and intercellular adhesion
%    - Does the mapping J/κ → p_0 hold quantitatively?
%    - The effective shape index p_eff measured from contours tests this directly

% 2. The quench experiment: what vertex models cannot show
%    - In vertex models, T1 transitions are instantaneous algorithmic events
%    - In our phase field, they are continuous, observable relaxation processes
%    - The quench reveals the energy landscape topology:
%      rigid phase = isolated minima with high barriers
%      fluid phase = flat landscape with connected minima
%    - Relaxation time divergence as a signature of criticality

% 3. Role of cell deformability (Saito & Ishihara 2024)
%    - Saito & Ishihara showed that cell deformability controls a fluid-to-fluid
%      transition distinct from the shape index transition
%    - Our model naturally includes deformability through the phase field profile
%    - The adhesion term modifies the interface width at cell-cell contacts
%    - This may lead to effects not captured by the vertex model

% 4. Comparison to other phase field adhesion studies
%    - Nonomura (2012): cortex overlap adhesion g(phi_i)g(phi_j)
%    - Najem & Grant (2016): range-field adhesion with auxiliary C_i
%    - Lober et al. (2015): bilinear overlap for repulsion; we flip sign for adhesion
%    - Wenzel & Voigt (2021): signed-distance repulsion, compared activity models
%    - Graham et al. (2024): active cell sorting, no adhesion term
%    - Key distinction: none mapped the rigidity transition via adhesion
%    - Our bilinear phi_i*phi_j is the simplest coupling; its longer range
%      vs interface-localized forms is a deliberate choice

% 5. Limitations
%    - 2D only (3D extension straightforward with ∫φiφj dV)
%    - Fixed κ — scanning J/κ by varying J at fixed κ
%    - No substrate friction, no cell division/death
%    - Overdamped dynamics (no inertia)

\todo{Write discussion after results are available.}


% ══════════════════════════════════════════════════════════════════════════════
%  VIII.  CONCLUSION
% ══════════════════════════════════════════════════════════════════════════════
\section{Conclusion}
\label{sec:conclusion}

\todo{%
Summarize: (1)~adhesion-induced rigidity transition exists in the phase field
model and is controlled by $\Jk$; (2)~the quench protocol reveals continuous
relaxation paths inaccessible to vertex models; (3)~the $(\Jk,\,\vA)$ phase
diagram is qualitatively (quantitatively?) analogous to the $(p_0,\,v_0)$
diagram; (4)~the effective shape index $\peff$ connects the two frameworks.
Outlook: 3D extension, differential adhesion for heterogeneous tissues,
Griffiths singularity under adhesion.}


% ══════════════════════════════════════════════════════════════════════════════
%  ACKNOWLEDGMENTS
% ══════════════════════════════════════════════════════════════════════════════
\begin{acknowledgments}
Computational resources were provided by the Digital Research Alliance of
Canada (Nibi cluster).
\todo{Add funding acknowledgments.}
\end{acknowledgments}

% ══════════════════════════════════════════════════════════════════════════════
%  BIBLIOGRAPHY
% ══════════════════════════════════════════════════════════════════════════════
\bibliography{references}

\end{document}
